```latex
\documentclass{article}
\usepackage{amsmath, amssymb}
\usepackage{enumitem}
\usepackage{geometry}
\geometry{margin=1in}

\begin{document}

\section*{Worked Solutions}

\subsection*{Problem 1: Basic Arithmetic}
\textbf{Question:} Calculate \( 15 \times (4 + 3) - 28 \div 7 \).

\textbf{Solution:}
\begin{align*}
15 \times (4 + 3) - 28 \div 7 &= 15 \times 7 - 28 \div 7 \quad \text{(Parentheses first)} \\
&= 105 - 4 \quad \text{(Multiplication and division from left to right)} \\
&= 101 \quad \text{(Subtraction)}
\end{align*}
\textbf{Answer:} \( \boxed{101} \)

\subsection*{Problem 2: Solving a Linear Equation}
\textbf{Question:} Solve for \( x \): \( 3x - 7 = 14 \).

\textbf{Solution:}
\begin{align*}
3x - 7 &= 14 \\
3x &= 14 + 7 \quad \text{(Add 7 to both sides)} \\
3x &= 21 \\
x &= \frac{21}{3} \quad \text{(Divide both sides by 3)} \\
x &= 7
\end{align*}
\textbf{Answer:} \( \boxed{x = 7} \)

\subsection*{Problem 3: Quadratic Equation}
\textbf{Question:} Solve \( x^2 - 5x + 6 = 0 \).

\textbf{Solution:}
We can factor the quadratic:
\[
x^2 - 5x + 6 = (x - 2)(x - 3) = 0
\]
Set each factor equal to zero:
\begin{align*}
x - 2 &= 0 \quad \Rightarrow \quad x = 2 \\
x - 3 &= 0 \quad \Rightarrow \quad x = 3
\end{align*}
\textbf{Answer:} \( \boxed{x = 2 \text{ or } x = 3} \)

\subsection*{Problem 4: Geometry - Area of a Circle}
\textbf{Question:} Find the area of a circle with radius \( r = 5 \) cm. Use \( \pi \approx 3.1416 \).

\textbf{Solution:}
The area formula is \( A = \pi r^2 \).
\[
A = \pi \times (5)^2 = \pi \times 25 \approx 3.1416 \times 25 = 78.54
\]
\textbf{Answer:} \( \boxed{78.54 \text{ cm}^2} \)

\subsection*{Problem 5: Word Problem}
\textbf{Question:} A store sells apples for \$2 each and oranges for \$3 each. If you buy 4 apples and 2 oranges, how much do you pay?

\textbf{Solution:}
Cost of apples: \( 4 \times 2 = 8 \) dollars. \\
Cost of oranges: \( 2 \times 3 = 6 \) dollars. \\
Total cost: \( 8 + 6 = 14 \) dollars.
\[
\text{Total} = (4 \times 2) + (2 \times 3) = 8 + 6 = 14
\]
\textbf{Answer:} \( \boxed{\$14} \)

\end{document}
```